\chapter{Circuits and Lower Bounds}
\label{ch:circuits}

A proof that $\Pc \neq \NP$ would most likely be a proof that some problem in $\NP$ requires Boolean circuits of superpolynomial size. A proof that $\Pc = \BPP$ would most likely go through a proof that some problem in exponential time requires circuits of exponential size. Both are circuit lower bounds, and this chapter surveys what is known about them: Shannon\index{Shannon, Claude}'s counting argument, the lower bounds for restricted circuit classes, the few results for general circuits, and the connection to derandomization that makes the subject central to the thesis.

\section{Shannon's counting argument}

Claude Shannon observed in 1949 that most Boolean functions need large circuits. There are $2^{2^n}$ functions of $n$ variables. A circuit with $s$ gates, each a binary gate on two earlier gates or inputs, can be described in $O(s \log s)$ bits, so there are at most $2^{O(s \log s)}$ circuits of size $s$. For these to compute all functions one needs $s \log s \ge 2^n$, so $s \ge 2^n / n$ roughly. Lupanov showed in 1958 that $2^n / n$ gates also suffice for every function, so the typical function has complexity about $2^n / n$ and the bound is tight.

The argument is nonconstructive in the strongest sense: it shows that hard functions exist without exhibiting one. The problem of circuit lower bounds is the problem of finding an explicit hard function, one in $\NP$ or in $\mathrm{E}$, and the gap between what counting gives for a random function and what anyone can prove for a specific one is the gap between $2^n / n$ and $3.1 n$.

\section{Restricted classes}

Progress has come by restricting the circuits. Three results of the 1980s are the landmarks.

Constant depth. The class $\mathrm{AC}^0$ consists of circuits of polynomial size and constant depth with unbounded fan in AND and OR gates. Furst, Saxe and Sipser\index{Sipser, Michael} in 1984 and Ajtai in 1983 showed that parity is not in $\mathrm{AC}^0$, and H{\aa}stad in 1986 sharpened the bound to size $2^{\Omega(n^{1/(d-1)})}$ for depth $d$. The technique is the switching lemma: a random restriction that fixes most of the variables turns a small depth $d$ circuit into a small depth $d - 1$ circuit, and iterating collapses the circuit to a constant while parity survives every restriction. The result separates $\mathrm{AC}^0$ from the classes above it and gives, through Nisan's work, pseudorandom generator\index{pseudorandom generator}s against $\mathrm{AC}^0$, one of the few settings in which a pseudorandom generator is known unconditionally.

Monotone circuits. Razborov\index{Razborov, Alexander} in 1985 showed that clique requires monotone circuits, those with AND and OR but no NOT, of superpolynomial size, and Alon and Boppana improved this to exponential. The technique is the approximation method: one shows that any small monotone circuit can be approximated gate by gate by a function of a special form, with a small error at each gate, and that the special form cannot approximate clique. This was the first superpolynomial lower bound for a natural $\NP$ complete problem against a natural circuit class, and it raised hopes that the monotone restriction could be removed. Razborov himself showed in 1985 that it could not, in the sense that perfect matching, which is in $\Pc$, also requires superpolynomial monotone circuits, so that monotone complexity does not track $\Pc$.

Modular gates. The class $\mathrm{ACC}^0$ adds gates that compute a sum modulo a fixed $m$. Razborov and Smolensky showed in 1987 that for $m$ a prime power, $\mathrm{AC}^0$ with $\mathrm{MOD}_m$ gates cannot compute $\mathrm{MOD}_q$ for a different prime $q$. For composite $m$ such as $6$ nothing was known until Williams\index{Williams, Ryan} in 2011 proved that $\mathrm{NEXP} \not\subseteq \mathrm{ACC}^0$, by an entirely new method that converts a slightly faster than brute force satisfiability algorithm for $\mathrm{ACC}^0$ circuits into a lower bound. Williams's result is the first lower bound in decades against a class not already handled by the 1980s techniques, and his method, algorithms to lower bounds, is a live hope for progress against general circuits.

\section{General circuits}

Against unrestricted circuits the record is dismal. The best lower bound for an explicit function in $\NP$ over the full binary basis is $3.1 n - o(n)$ by Li and Yang~\cite{LiYang} in 2022, improving $(3 + 1/86) n - o(n)$ by Find, Golovnev, Hirsch and Kulikov in 2016 and improving Blum's $3n$ from 1984, and over the basis without parity gates it is $5n - o(n)$ by Iwama and Morizumi; for formulas, circuits with fan out one, the best is $n^{3 - o(1)}$ by H{\aa}stad, improving Andreev. No superlinear bound is known for circuits and no superpolynomial bound is known for general circuits or formulas. The gap between the linear bounds one can prove and the $2^n / n$ one knows to exist is the whole content of the $\Pc$ versus $\NP$ problem seen from the circuit side.

The reason is the barriers of Chapter~\ref{ch:barriers}. The techniques that work for restricted classes, random restrictions and approximation, are natural in the sense of Razborov and Rudich: they establish that the hard function has a property that is easy to check and shared by many functions. Such techniques cannot prove lower bounds against general circuits unless pseudorandom functions do not exist, and pseudorandom functions exist if one way functions do. So the methods of the 1980s are known not to extend, and the method of Williams, which is not natural, has so far reached only $\mathrm{ACC}^0$.

\section{Hardness versus randomness again}

The derandomization theorems of Chapter~\ref{ch:pseudo} need hardness of a specific kind: a function in $\mathrm{E}$ that requires circuits of size $2^{\Omega(n)}$. By Shannon's argument such functions exist in abundance; by the time hierarchy theorem $\mathrm{E}$ contains functions of every complexity up to its own; but no explicit function in $\mathrm{E}$ is known to be hard, for the same reason that none in $\NP$ is. The hypothesis of Impagliazzo\index{Impagliazzo and Wigderson}\index{Wigderson, Avi} and Wigderson is therefore of the same character as $\Pc \neq \NP$, a circuit lower bound for an explicit function, and it is believed for the same reason: counting shows hard functions are generic, and there is no known reason that explicit functions should be special.

There is a partial converse that sharpens the picture. Impagliazzo, Kabanets\index{Kabanets and Impagliazzo} and Wigderson showed in 2002 that if $\mathrm{NEXP} \subseteq \Ppoly$ then $\mathrm{NEXP} = \mathrm{MA}$, the class of problems with a randomized verifier, and Kabanets and Impagliazzo showed that derandomizing polynomial identity testing implies either $\mathrm{NEXP} \not\subseteq \Ppoly$ or a lower bound for the permanent. Derandomization and lower bounds are thus two faces of one coin, and a proof of $\Pc = \BPP$ by any route would carry a lower bound with it.

\section{What this means for the thesis}

The thesis of this book is stated on the quantile axis and says nothing directly about circuits. But any resolution of $\Pc = \RP$ along the lines of Chapter~\ref{ch:prp}, a deterministic polynomial size set that hits every dense level set of a polynomial time computable function, is a hitting set\index{hitting set} generator, and a hitting set generator implies a circuit lower bound by the Kabanets Impagliazzo route. The Monte Carlo program cannot evade this. If it succeeds in derandomizing the single draw of the $\RP$ walk, it will have proved that some explicit function needs large circuits, and the proof will have to be non natural and non relativiz\index{relativization}ing. We do not know how to do this. We say where the draw must land; we do not say how to land it without a coin.
